\section{主观要件 - 事实认识错误}
\begin{enumerate}
  \item \kd{事实认知错误} = 主客观一致的反面。\zd{核心。要做到主客观相一致}。
  \item 体系审查。\zd{\blx{先客观，后主观}}。
        \begin{enumerate}
          \item 客观。
                \begin{enumerate}
                  \item 危害行为 √ 打 ×。
                        \begin{enumerate}
                          \item 打 ×。\al{沙漠稻草人 \yjt 对象不能犯}。
                          \item 打 √ \xjt。
                        \end{enumerate}
                  \item 危害结果。
                \end{enumerate}
          \item 主观。有无有犯罪故意。
                \begin{enumerate}
                  \item 没有。\zd{缺乏犯罪故意型认识错误}。\al{甲以为前方是野兔，结果打死乙 \yjt 过失还是意外事件}。
                  \item 有。\zd{具有犯罪故意型认识错误}。\al{甲以为前方是仇人乙，结果打死丙/打到文物 \yjt 能否构成既遂}。
                        \begin{enumerate}
                          \item 具体的错误/同一犯罪构成内的错误。整体行为还是评价成故意杀人，你想杀人，杀的也是个人。
                        \end{enumerate}

                        \begin{kk}{对象错误}
                          \al{甲\xhx{误以为}路人是仇人丙开枪打死 \yjt 故意杀人既遂}。
                          \begin{enumerate}
                            \item 人是抽象的，不要求真正是丙，\zd{身份不重要}。只是杀人\zd{动机没实现}。
                            \item \zd{动机是动机，故意是故意}。
                          \end{enumerate}
                        \end{kk}

                        \begin{kk}{打击错误。三角关系。预定目标、实害对象}
                          \al{甲杀乙，错杀路过的丙}。
                          \begin{enumerate}
                            \item 对实害结果不能是故意（直接、间接均可），\zd{只能是过失}（如果主观认识到可能会把丙弄死，那就主客观一致了）。
                            \item \zd{具体符合说}。尊重案件事实，分别分析。\zd{预定未遂，实害过失，想象竞合择一重 \yjt 未遂}。
                            \item 法定符合说。为了严惩凶手，\zd{过失拟制为故意，预定未遂，实害既遂，想象竞合择一重 \yjt 既遂}。\al{甲开枪杀乙，子弹穿过乙还杀了丙}。
                            \item 和打击错误结合。\zd{位置比身份重要，把姓名抹掉}。\al{二皇子刺杀太子案}。
                            \item 和偶然防卫的结合。\al{猪队友帮倒忙。甲乙联合枪杀丙，甲射偏射打死了乙，丙免遭死亡}。
                          \end{enumerate}
                        \end{kk}

                        \begin{kk}{因果关系错误。设想的因果流程不一样}
                          \begin{enumerate}
                            \item 狭义的因果关系错误。\zd{只有一个行为}。\al{想毒死乙却被噎死 \yjt 没有因果关系 \yjt 未遂、甲相推乙入井淹死却摔死 \yjt 有 \yjt 既遂}。
                          \end{enumerate}
                          \tcblower
                          \zd{事前故意/结果的推迟发生}。\zd{前行为异常（昏迷）+ 后行为}。\al{甲杀乙打的重伤\xhx{昏迷}以为乙死亡，抛尸水里}。
                          \begin{enumerate}
                            \item \zd{多数说}。二因一果 \yjt 前杀人既遂 \yjt 后行为过失致人死亡 \yjt 吸收 \yjt \zd{定杀人既遂}。
                            \item 少数说。后因一果 \yjt 前杀人未遂 \yjt \zd{后行为过失致人死亡} \yjt 并罚 \yjt \zd{定杀人未遂 + 过失致人死亡}。
                            \item 前行为不异常 + 后行为。正常处理。前杀人未遂，后过失致人死亡，并罚。
                          \end{enumerate}
                          \zd{结果的提前实现}。前行为（预备行为）+ 计划行为。\zd{预备行为转为了实行行为}（推定的）。
                          \begin{enumerate}
                            \item 关键。判断\zd{实害结果发生在哪个阶段}。
                            \item \al{妻子计划第二天毒死丈夫，丈夫提前回来} \yjt 放毒酒时\zd{未产生}现实紧迫危险。
                            \item 未着手。故意杀人犯罪预备、过失致人死亡（安全注意义务），想象竞合择一重。
                            \item \al{先安眠药弄晕再勒死，结果安眠药过量} \yjt 安眠药\zd{有现实紧迫危险}。
                            \item 投放时已着手。故意杀人既遂。
                            \item \al{打算自杀浇满汽油，打算吸根烟再死，结果点烟着火} \yjt \zd{已产生}现实紧迫危险。
                          \end{enumerate}
                        \end{kk}

                        \begin{enumerate}
                          \item 抽象的错误/不同犯罪构成间的错误。预定目标和实害对象\zd{法益性质种类不同}。\zd{分别分析法}。
                          \item \al{甲想打碎珍贵蜡像，结果打偏打死了旁边的乙}。预定毁坏财物未遂，实害过失致人死亡，想象竞合择一重。
                          \item 找俩包容评价关系的罪名结合考。\al{甲在车站行窃时盗得一提包，回家一看才发现包内仅有一把手枪。怕被人发现，藏在浴缸下}。
                        \end{enumerate}
                \end{enumerate}
        \end{enumerate}
        \begin{kk}{对象错误和打击错误的区分。两步走}
          \begin{enumerate}
            \item 对实害对象的实害结果 \yjt 过失 \hz 意外事件 \yjt 打击错误。
            \item 故意 \yjt 对实害对象的身份有/没有认识错误 \yjt 对象错误/普通的间接故意。
          \end{enumerate}
        \end{kk}
\end{enumerate}
\centering
\begin{forest}
  my tree,
  for tree={
    grow=south,
  },
  [体系审查
    [客观
      [危害行为
        [没有
          [沙漠稻草人
            [对象不能犯]
          ]
        ]
        [有]
      ]
      [危害结果]
    ]
    [主观\\（有无有犯罪故意）,align=center,base=bottom
      [没有\\（缺乏犯罪故意型认识错误）,align=center,base=bottom
        [事实认识错误
          [焦点\\过失 \hz 意外,align=center,base=bottom]
        ]
        [法律认识错误
          [违法性认识可能性]
        ]
      ]
      [有\\（具有犯罪故意型认识错误）,align=center,base=bottom
        [预定目标和实害对象法益性质种类
          [具体的错误\\（同一犯罪构成内的错误）,align=center,base=bottom
              [故意
                [对象认识错误
                  [有\\（对象错误）,align=center,base=bottom]
                  [没有\\（普通的间接故意）,align=center,base=bottom]
                ]
              ]
              [过失
                [打击错误
                  [具体符合说]
                  [法定符合说]
                ]
              ]
              [因果关系
                [狭义的]
                [结果的推迟发生]
                [结果的提前实现]
              ]
          ]
          [抽象的错误\\（不同犯罪构成间的错误）,align=center,base=bottom
            [分别分析]
          ]
        ]
      ]
    ]
  ]
\end{forest}


